% Generated draft: requires mathematical and source review.
\documentclass[11pt]{article}
\usepackage[margin=25mm]{geometry}
\usepackage{fontspec}
\setmainfont{Latin Modern Roman}
\usepackage{amsmath,amssymb,mathtools,unicode-math}
\setmathfont{Latin Modern Math}
\usepackage{enumitem,needspace,xurl,hyperref}
\setlength{\parindent}{0pt}
\setlength{\parskip}{5pt}
\setlength{\emergencystretch}{3em}
\newcommand{\N}{\mathbb{N}}
\newcommand{\Z}{\mathbb{Z}}
\newcommand{\Q}{\mathbb{Q}}
\newcommand{\R}{\mathbb{R}}
\newcommand{\C}{\mathbb{C}}
\newcommand{\F}{\mathbb{F}}
\newcommand{\A}{\mathbb{A}}
\newcommand{\PP}{\mathbb{P}}
\newcommand{\E}{\mathbb{E}}
\newcommand{\eps}{\varepsilon}
\newcommand{\dd}{\,\mathrm{d}}
\begin{document}
\section*{UnsolvedMath ID: 4400005}
\textbf{Source problem code:} catalogue entry
\subsection{Definitions and mathematical statement}
Let $(M,h)$ be a compact Riemannian manifold of constant negative sectional curvature. Write
\[
 SM=\{(x,v)\in TM:h_x(v,v)=1\}
\]
for its unit tangent bundle and $g_t:SM\to SM$ for its unit-speed geodesic flow, $t\in\R$. For a Borel probability measure $\mu$ on $SM$, invariance under $g_t$ means $(g_t)_*\mu=\mu$.

The supplied question asks whether there exists a Borel probability measure $\mu$ satisfying
\[
 (g_1)_*\mu=\mu,
 \qquad\text{but}\qquad
 \exists t\in\R\quad(g_t)_*\mu\ne\mu.
\]
Equivalently, the second condition asks for some time $t$ and Borel set $A\subseteq SM$ with $\mu(g_t^{-1}A)\ne\mu(A)$. The state space is the unit tangent bundle on which the geodesic flow acts, not just the base manifold.

The source begins with an arbitrary $M$, but does not clarify whether the existence question is required for every such fixed manifold or only for at least one example. It gives no curvature normalization, although rescaling the metric changes the relation between geodesic periods and the specified time one. It imposes no non-atomicity, full-support, entropy, regularity, ergodicity, or exclusion of periodic-orbit-supported measures. Those additional restrictions must not be inferred from the difficulty label or the discussion of rigidity results. In particular possible periodic-orbit cases and the role of time normalization require source review before interpreting the broad formulation as a specific unresolved rigidity problem.
\subsection{Short English statement}
Can a probability measure on the unit tangent bundle of a compact constant-negatively-curved manifold be invariant under one unit of geodesic time without being invariant under every real time? The supplied summary reports rigidity results under additional entropy or regularity hypotheses, and says it did not verify a resolution of the fully general existence question. It gives no precise theorem or primary reference for that progress.

The registry marks partial progress. The 2010 source's open label and the later summary are supplied assessments, with no later authoritative resolution check recorded; they do not add the restrictions absent from the statement.
\subsection{Sources}
\begin{itemize}
\item Local UnsolvedMath registry, permanent ID 4400005, source code catalogue entry, ``Pingree open problems --- Ledrappier problem 1,'' proposed by François Ledrappier, 2010. Dataset curation and metadata: Ulam AI and UnsolvedMath Contributors, CC BY 4.0.
\item \emph{Open problems from the 3rd Pingree Workshop on Dynamical Systems} (2010), François Ledrappier, Problem 1, \url{https://math.huji.ac.il/~mhochman/open-problems/pingree-open-problems.pdf}; pdf-text extraction, accessed 2026-07-29.
\item The supplied summary refers to strong rigidity results under entropy or regularity assumptions without identifying papers or exact hypotheses. No additional citation is inferred.
\item Source provenance marks status evidence NEEDS\_REVIEW and rights NEEDS\_REVIEW: the source presents the item as open, no later authoritative resolution check is recorded, and redistribution permission was not established despite verified public access. This qualification is separate from the dataset licence.
\end{itemize}
\end{document}
